\section{Related Rates}\label{sec:related_rates}

When two quantities are related by an equation, knowing the value of one quantity can determine the value of the other. For instance, the circumference and radius of a circle are related by $C=2\pi r$; knowing that $C = 6\pi$ in determines the radius must be 3 in.

The topic of \textbf{related rates} takes this one step further: knowing the \emph{rate} at which one quantity is changing can determine the \emph{rate} at which the other changes.\index{related rates}

\youtubeVideo{tZl5h7590go}{Related Rates \#6 --- Rate at Which the Circumference of a Circle is Changing}

\mnote{\textbf{Note:} This section relies heavily on implicit differentiation, so referring back to \autoref{sec:imp_deriv} may help.}

We demonstrate the concepts of related rates through examples.

\begin{example}[Understanding related rates]\label{ex_rr1}%
The radius of a circle is growing at a rate of 5 in/hr. At what rate is the circumference growing?
\solution
The circumference and radius of a circle are related by $C = 2\pi r$. We are given information about how the length of $r$ changes with respect to time; that is, we are told $\frac{\dd r}{\dd t} = 5$ in/hr. We want to know how the length of $C$ changes with respect to time, i.e., we want to know $\frac{\dd C}{\dd t}$. 

Implicitly differentiate both sides of $C = 2\pi r$ with respect to $t$:
\begin{align*}
C 	&= 2\pi r\\
\frac{\dd}{\dd t}\bigl(C\bigr) &= \frac{\dd}{\dd t}\bigl(2\pi r\bigr) \\
\frac{\dd C}{\dd t} &= 2\pi\frac{\dd r}{\dd t}.
\end{align*}
As we know $\frac{\dd r}{\dd t} = 5$ in/hr, we know
\[
 \frac{\dd C}{\dd t} = 2\pi 5 = 10\pi %\approx 31.4
 \text{ in/hr.}
\]
\end{example}

Before we look at another example, we'll state a few ideas on approaching these problems.

\begin{keyidea}[Solving Related Rates Problems]\label{idea:related_rates}%
~\\[-2\baselineskip]
\parbox[t]{\linewidth}{\begin{enumerate}
	\item	Understand the problem. Clearly identify what quantities are given and what are to be found. Make a sketch if helpful.
	\item	Assign mathematical notation to all quantities, including those that are functions of time.
	\item	Create an equation relevant to the context of the problem, using the information given.
	\item	Substitute constant quantities and if necessary, use the given information to eliminate other variables.
	\item	Use the Chain Rule to differentiate both sides of the equation.
	\item	Substitute the known quantities, and solve for the unknown rate.
\end{enumerate}}
\end{keyidea}

The important thing to remember is that you must differentiate before you substitute varying values.  Otherwise, you'll substitute a constant for what should be a variable, and it's derivative will be zero.  Consider another example.

\begin{example}[Finding related rates]\label{ex_rr2}%
Water streams out of a faucet at a rate of $2$ in$^3$/s onto a flat surface at a constant rate, forming a circular puddle that is $1/8$ in deep. 
\begin{enumerate}
	\item	At what rate is the area of the puddle growing?
	\item	At what rate is the radius of the circle growing?
\end{enumerate}
\solution
\begin{enumerate}
\item We can answer this question two ways: using ``common sense'' or related rates. The common sense method states that the volume of the puddle is growing by $2$ in$^3$/s, where
\[\text{volume of puddle } = \text{ area of circle } \times \text{ depth.}\]
Since the depth is constant at $1/8$ in, the area must be growing by 16in$^2$/s.

This approach reveals the underlying related-rates principle. Let $V$ and $A$ represent the Volume and Area of the puddle. We know $V= A\times \frac18$. Take the derivative of both sides with respect to $t$, employing implicit differentiation.
\begin{align*}
V &= \frac18A\\
\frac{\dd}{\dd t}\bigl(V\bigr) &= \frac{\dd}{\dd t}\left(\frac18A\right)\\
\frac{\dd V}{\dd t} &= \frac18\frac{\dd A}{\dd t}
\end{align*} 
As $\frac{\dd V}{\dd t} = 2$, we know $2 = \frac18\frac{\dd A}{\dd t}$, and hence $\frac{\dd A}{\dd t} = 16$. Thus the area is growing by 16 in$^2$/s.

\item		To start, we need an equation that relates what we know to the radius.
We know that $V=\pi r^2h=\frac\pi8r^2$.  Implicitly derive both sides with respect to $t$:
\begin{align*}
 V &= \frac\pi8 r^2 \\
 \frac{\dd}{\dd t}\bigl(V\bigr) &= \frac{\dd}{\dd t}\bigl(\frac\pi8 r^2\bigr) \\
 \frac{\dd V}{\dd t} &= \frac\pi8\cdot2r\frac{\dd r}{\dd t}=\frac\pi4 r\frac{\dd r}{\dd t}
\end{align*}
Solving for $\dfrac{\dd r}{\dd t}$, we have
$\ds
\frac{\dd r}{\dd t}=\frac{\frac{\dd V}{\dd t}}{\frac\pi4 r}
=\frac{2}{\frac\pi4 r}=\frac8{\pi r}.
$
% We just learned something about the surface area of the circular puddle, and we know $A = \pi r^2$. We should be able to learn about the rate at which the radius is growing with this information. 
%
%Implicitly derive both sides of $A=\pi r^2$ with respect to $t$:
%\begin{align*}
%	A 	&= \pi r^2 \\
%	\frac{\dd}{\dd t}\bigl(A\bigr) &= \frac{\dd}{\dd t}\bigl(\pi r^2\bigr)\\
%	\frac{\dd A}{\dd t} &= 2\pi r\frac{\dd r}{\dd t}
%\end{align*}
%
%Our work above told us that $\frac{\dd A}{\dd t} = 16$ in$^2$/s. Solving for $\frac{\dd r}{\dd t}$, we have \[\frac{\dd r}{\dd t} = \frac{8}{\pi r}.\]

Note how our answer is not a number, but rather a function of $r$. In other words, \emph{the rate at which the radius is growing depends on how big the circle already is.} If the circle is very large, adding 2 in$^3$ of water will not make the circle much bigger at all. If the circle is dime-sized, adding the same amount of water will make a radical change in the radius of the circle.

In some ways, our problem was (intentionally) ill-posed. We need to specify a current radius in order to know a rate of change. When the puddle has a radius of 10 in, the radius is growing at a rate of
\[
\frac{\dd r}{\dd t} = \frac{8}{10\pi} = \frac{4}{5\pi} \approx 0.25\text{ in/s}.
\]
\end{enumerate}
\end{example}

\begin{example}[Studying related rates]\label{ex_rr3}%
Radar guns measure the rate of distance change between the gun and the object it is measuring. For instance, a reading of ``55 mph'' means the object is moving away from the gun at a rate of 55 miles per hour, whereas a measurement of ``$-25$ mph'' would mean that the object is approaching the gun at a rate of 25 miles per hour.

If the radar gun is moving (say, attached to a police car) then radar readouts are only immediately understandable if the gun and the object are moving along the same line. If a police officer is traveling 60 mph and gets a readout of 15 mph, he knows that the car ahead of him is moving away at a rate of 15 miles an hour, meaning the car is traveling 75 mph. (This straight-line principle is one reason officers park on the side of the highway and try to shoot straight back down the road. It gives the most accurate reading.)

\mtable{A sketch of an officer attempting to measure the speed of a car in \autoref{ex_rr3}.}{fig:rr3}{\begin{tikzpicture}[alt={Red-dot officer at origin; vertical axis points north. Car sits 0.5 unit east on road 0.5 unit north; diagonal C joins them, car arrow points east.},>=latex]
\draw(0,0) -- (3,0) node [pos=.5,above] {\scriptsize $B=1/2$} -- (0,-3) node [pos=.4,below] {\scriptsize $C$}-- (0,0) node [pos=.5,rotate=90,shift={(.1,.2)}] {\scriptsize $A=1/2$};
\draw[thick] (0,-3.5) -- (0,.5) node [above] {\scriptsize $N$};
\draw[thick] (-.5,0) -- (3.5,0) node [right] {\scriptsize $E$};
\filldraw[fill={\colortwo},draw={\colortwo}] (0,-3) circle (2pt) node [right] {\scriptsize Officer};
\filldraw[fill=black,draw=black] (3,0) circle (2pt) node [below ] {\scriptsize \quad Car};
\draw[->,thick] (-.2,-3.3) -- (-.2,-2.7);
\draw[->,thick] (2.7,.2) -- (3.3,.2);
\end{tikzpicture}}

Suppose an officer is driving due north at 30 mph and sees a car moving due east, as shown in \autoref{fig:rr3}. Using his radar gun, he measures a reading of 20 mph. By using landmarks, he believes both he and the other car are about 1/2 mile from the intersection of their two roads. 

If the speed limit on the other road is 55 mph, is the other driver speeding?
\solution
Using the diagram in \autoref{fig:rr3}, let's label what we know about the situation. As both the police officer and other driver are 1/2 mile from the intersection, we have $A = 1/2$, $B = 1/2$, and through the Pythagorean Theorem, $C = 1/\sqrt{2}\approx 0.707$. 

We know the police officer is traveling at 30 mph; that is, $\frac{\dd A}{\dd t} = -30$. The reason this rate of change is negative is that $A$ is getting smaller; the distance between the officer and the intersection is shrinking. The radar measurement is $\frac{\dd C}{\dd t} = 20$. We want to find $\frac{\dd B}{\dd t}$. 

We need an equation that relates $B$ to $A$ and/or $C$. The Pythagorean Theorem is a good choice: $A^2+B^2 = C^2$. Differentiate both sides with respect to $t$:
\begin{align*}
	A^2 + B^2 &= C^2 \\
	\frac{\dd}{\dd t}\left(A^2+B^2\right) &= \frac{\dd}{\dd t}\left(C^2\right) \\
	2A\frac{\dd A}{\dd t} + 2B\frac{\dd B}{\dd t} &= 2C\frac{\dd C}{\dd t}
\end{align*}

\mnote{\textbf{Note:} \autoref{ex_rr3} is both interesting and impractical. It highlights the difficulty in using radar in a non-linear fashion, and explains why ``in real life'' the police officer would follow the other driver to determine their speed, and not pull out pencil and paper.\bigskip\\
The principles here are important, though. Many automated vehicles make judgments about other moving objects based on perceived distances, radar-like measurements and the concepts of related rates.}

We have values for everything except $\frac{\dd B}{\dd t}$. Solving for this we have 
\[
\frac{\dd B}{\dd t} = \frac{C\frac{\dd C}{\dd t}- A\frac{\dd A}{\dd t}}{B} \approx 58.28\text{ mph}.
\]
		
The other driver appears to be speeding slightly.	
\end{example}

\begin{example}[Studying related rates]\label{ex_rr4}%
A camera is placed on a tripod 10 ft from the side of a road. The camera is to turn to track a car that is to drive by at 100 mph for a promotional video. The video's planners want to know what kind of motor the tripod should be equipped with in order to properly track the car as it passes by.  \autoref{fig:rr4} shows the proposed setup.

\mtable{Tracking a speeding car (at left) with a rotating camera.}{fig:rr4}{\begin{tikzpicture}[alt={Red car on road, arrow east 100 mph; camera 10 ft below road. Horizontal gap x, diagonal to camera makes angle θ.}
,>=latex]
\draw (-3,0) -- (0,0) -- (0,-2) node [shift={(-5pt,10pt)}] {\small $\theta$} node [right,pos=.5] {\scriptsize 10ft}-- cycle;
\draw [<->, thick] (-3.5,0) -- node [below,pos=.5] {\scriptsize $x$} (.5,0);
\draw [fill={\colortwo},draw={\colortwo}] (-3,0) circle (2pt);
\draw [fill=black] (0,-2) circle (2pt);
\draw [->] (-3,.2) -- (-2,.2) node [above,pos=.5] {\scriptsize 100mph};
\end{tikzpicture}}

How fast must the camera be able to turn to track the car?
\solution
We seek information about how fast the camera is to \emph{turn}; therefore, we need an equation that will relate an angle $\theta$ to the position of the camera and the speed and position of the car.

\autoref{fig:rr4} suggests we use a trigonometric equation. Letting $x$ represent the distance the car is from the point on the road directly in front of the camera, we have
\begin{equation}\tan \theta = \frac{x}{10}.\label{eq:rr4}\end{equation}%
As the car is moving at 100 mph, we have $\frac{\dd x}{\dd t} = -100$ mph (as in the last example, since $x$ is getting smaller as the car travels, $\frac{\dd x}{\dd t}$ is negative). We need to convert the measurements so they use the same units; rewrite -100 mph in terms of ft/s:
\[
\frac{\dd x}{\dd t} = -100\frac{\text{m}}{\text{hr}} = -100\frac{\text{m}}{\text{hr}}\cdot5280\frac{\text{ft}}{\text{m}}\cdot\frac{1}{3600}\frac{\text{hr}}{\text{s}} =-146.\overline{6}\text{ ft/s}.
\]
Now take the derivative of both sides of \autoeqref{eq:rr4} using implicit differentiation:
\begin{align}
	\tan \theta &= \frac{x}{10} \notag \\
	\frac{\dd}{\dd t}\bigl(\tan \theta\bigr)
	&= \frac{\dd}{\dd t}\left(\frac{x}{10}\right)\notag \\
	\sec^2\theta\,\frac{\dd\theta}{\dd t} &= \frac{1}{10}\frac{\dd x}{\dd t}\notag\\
	\frac{\dd\theta}{\dd t} &= \frac{\cos^2\theta}{10}\frac{\dd x}{\dd t}\label{eq:rr4b}
\end{align}
We want to know the fastest the camera has to turn. Common sense tells us this is when the car is directly in front of the camera (i.e., when $\theta = 0$). Our mathematics bears this out. In \autoeqref{eq:rr4b} we see this is when $\cos^2\theta$ is largest; this is when $\cos \theta = 1$, or when $\theta = 0$.

With $\frac{\dd x}{\dd t} \approx -146.67$ ft/s, we have
\[
 \frac{\dd\theta}{\dd t}
 = -\frac{1\text{ rad}}{10\text{ ft}}146.67\text{ ft/s}
 = -14.667\text{ radians/s}.
\]
We find that $\frac{\dd\theta}{\dd t}$ is negative; this matches our diagram in \autoref{fig:rr4} for $\theta$ is getting smaller as the car approaches the camera.	
	
What is the practical meaning of $-14.667$ radians/s? Recall that 1 circular revolution goes through $2\pi$ radians, thus $14.667$ rad/s means $14.667/(2\pi)\approx 2.33$ revolutions per second. The negative sign indicates the camera is rotating in a clockwise fashion.
\end{example}

% todo write an example with similar triangles

We introduced the derivative as a function that gives the slopes of tangent lines of functions. This 
%chapter emphasizes using the derivative in other ways. Newton's Method uses the derivative to approximate roots of functions; this
section stresses the ``rate of change'' aspect of the derivative to find a relationship between the rates of change of two related quantities. 
In the next section we use Extreme Value concepts to \emph{optimize} quantities. 

% todo write a problem with water in & water out
% todo write a problem with an isosceles triangle
% todo write a problem with a person, light & shadow
% todo write a problem with a television camera & rocket?
% todo write a problem with a lighthouse beam?

\printexercises{exercises/04-02-exercises}
